\documentclass[conference]{IEEEtran}
\usepackage[utf8]{inputenc}
\usepackage[T1]{fontenc}
\usepackage[spanish,es-tabla]{babel}
\usepackage{graphicx}
\usepackage{booktabs}
\usepackage{array}
\usepackage{url}

% Código en línea que admite guiones bajos y permite cortar rutas largas en tablas
\makeatletter
\g@addto@macro{\UrlBreaks}{\do\_\do\:}
\makeatother
\urlstyle{tt}
\newcommand{\cd}[1]{\path{#1}}

\begin{document}

\title{Diseño e Implementación de la API REST de la Plataforma Web TasteBox\\
(Caso 6 -- Proyecto Bases de Datos, Segundo Corte)}

\author{\IEEEauthorblockN{Camilo Arbey Cuadros Mendez, Juan Diego Sierra}
\IEEEauthorblockA{Ingeniería Multimedia\\
Universidad de San Buenaventura\\
Bogotá, Colombia}}

\maketitle

\begin{abstract}
Este documento extiende el análisis de requerimientos de TasteBox (corte 1) con el diseño de la capa de datos y de la API. Se formaliza el modelo entidad--relación y se argumenta el proceso de normalización hasta la forma normal de Boyce--Codd; se mapean las entidades a recursos REST y se justifican las decisiones de anidamiento, referencias y endpoints compuestos; y se describe la arquitectura por capas (rutas, controladores, servicios y repositorios) conectada a MySQL. La implementación en Node.js y Express cubre los diecinueve requerimientos funcionales, se documenta con OpenAPI y Swagger UI y se verifica con una colección de pruebas con casos de éxito y de error por endpoint.
\end{abstract}

\begin{IEEEkeywords}
TasteBox, API REST, modelo entidad--relación, normalización, MySQL, Express, OpenAPI.
\end{IEEEkeywords}

% La numeración continúa a partir de la sección XII del informe del corte 1
\setcounter{section}{12}

\input{secciones_corte2}

\begin{thebibliography}{9}
\bibitem{owasp} OWASP Foundation, ``Password Storage Cheat Sheet,'' OWASP Cheat Sheet Series, 2026. [Online]. Available: https://cheatsheetseries.owasp.org/cheatsheets/Password\_Storage\_Cheat\_Sheet.html
\bibitem{nist} National Institute of Standards and Technology, ``Digital Identity Guidelines: Authentication and Lifecycle Management,'' NIST Special Publication 800-63B. [Online]. Available: https://pages.nist.gov/800-63-3/sp800-63b.html
\bibitem{fielding} R. T. Fielding, ``Architectural Styles and the Design of Network-based Software Architectures,'' Ph.D. dissertation, Dept. Inf. Comput. Sci., Univ. California, Irvine, CA, USA, 2000.
\bibitem{rfc9110} R. Fielding, M. Nottingham, and J. Reschke, ``HTTP Semantics,'' IETF, RFC 9110, Jun. 2022.
\bibitem{openapi} OpenAPI Initiative, ``OpenAPI Specification v3.0.3,'' 2020. [Online]. Available: https://spec.openapis.org/oas/v3.0.3
\bibitem{rfc7519} M. Jones, J. Bradley, and N. Sakimura, ``JSON Web Token (JWT),'' IETF, RFC 7519, May 2015.
\end{thebibliography}

\end{document}
